% 第 38 章　追着光跑会看见什么
\chapter{追着光跑会看见什么}

据说，爱因斯坦十六岁那年在阿劳中学读书时，脑子里钻进了一个甩不掉的问题：如果我跑得和光一样快，我会看见什么？一列光波，会不会像一列被我追平的火车那样，在我身边“停下来”？这个问题他追了十年。十年后，他给出的答案没有把光解释成一种奇怪的东西，而是把我们从小相信的“速度相加”这条常识，连同时间和空间本身，一起请下了神坛。这一章，我们重走这条路。

\step{反常识开场}

先从一个你闭着眼都能答对的问题开始。

站台上，一列火车以 \(30\,\mathrm{m/s}\)（大约 \(108\,\mathrm{km/h}\)）从你面前驶过。车厢里有个人朝车头方向抛出一个球，球相对车厢的速度是 \(10\,\mathrm{m/s}\)。请问：在你看来，球的速度是多少？

\begin{center}
\begin{tikzpicture}[font=\small,>={Stealth}]
  \draw[thick] (-0.5,0) -- (10,0);
  \draw[fill=gray!10] (1,0.28) rectangle (6.6,1.5);
  \draw (1.7,0.28) circle (0.18);
  \draw (5.9,0.28) circle (0.18);
  \node at (4.9,1.05) {车厢};
  \fill (2.1,1.05) circle (0.07);
  \draw[-{Stealth},thick] (2.1,1.05) -- (3.7,1.05)
    node[midway,above] {球相对车 \(10\,\mathrm{m/s}\)};
  \draw[-{Stealth},thick] (1,-0.5) -- (4.6,-0.5)
    node[midway,below] {车相对地面 \(30\,\mathrm{m/s}\)};
  \node at (8.8,0.5) {站台上的你};
  \fill (8.8,0.15) circle (0.07);
\end{tikzpicture}
\end{center}

你大概想都不想：\(30+10=40\,\mathrm{m/s}\)。从小到大，这条规则从没骗过你：在自动扶梯上往前走，你相对地面的速度是扶梯速度加步行速度；顺水划船，岸上看到的船速是水速加划船速度；高速公路上两辆车对开，各自 \(100\,\mathrm{km/h}\)，迎面而来的车在你眼里以 \(200\,\mathrm{km/h}\) 冲过来。速度相加，天经地义。

现在，把那个抛球的人换成一个按亮手电筒的人。火车还是以 \(30\,\mathrm{m/s}\) 行驶，手电朝车头方向射出一束光。你在站台上用一台足够精密的仪器去测这束光的速度，会得到多少？

如果世界遵守“速度相加”，答案应该是光速加上车速：\(c+30\,\mathrm{m/s}\)。

但实验的答案——从十九世纪末到今天，无数次越来越精密的测量的答案——是不多不少，正好 \(c\)，大约 \(299\,792\,458\,\mathrm{m/s}\)。光源在动也好，在飞也好，测出来的光速永远是同一个数。

更刺激的还在后面。让两辆车头对头地开，各自 \(100\,\mathrm{km/h}\)，你坐在其中一辆里，会看到另一辆以 \(200\,\mathrm{km/h}\) 扑面而来。现在让两束光迎面飞行，你“坐”在其中一束光上（先别管怎么坐上去），另一束光相对你的速度是多少？不是 \(2c\)。还是 \(c\)。

这一章剩下的全部内容，就是弄明白：为什么球的速度要相加，光的速度偏偏不加；以及这个“偏偏”背后，自然界到底在告诉我们什么。

\step{先猜一猜}

在往下读之前，请拿出一张纸，认真写下你对下面三个问题的猜测。不要跳过去——本章末尾你会回来给自己的直觉打分。

第一题：火车上的手电朝前发光，站台上的仪器测到的光速，是比 \(c\) 大、等于 \(c\)，还是取决于火车速度？

第二题：假如你坐上一艘能以光速的 \(99\%\) 飞行的飞船，去追一束从地球发出的激光。在你看来，这束激光以多快的速度离你远去？是 \(0.01\,c\) 吗？

第三题：两束光在真空中迎面相遇。从其中一束光的“视角”看，另一束光的速度是多少？\(2c\)？

写好了吗？如果你三题都按“速度相加”来答，恭喜，你和公元一九〇〇年之前几乎所有的物理学家站在同一边——包括那些最聪明的大脑。接下来你会看到，问题恰恰出在这条从没人怀疑过的“常识”上。

还要提醒一句：本章的怪事不在于“光很特殊”。读完你会发现，真正被修改的是“速度”这个词本身的运算规则，而光只是第一个把这件事暴露出来的信使。

\step{动手实验}

\begin{experiment}[用微波炉和一盘巧克力称出光速]
光速接近每秒三十万公里，听起来是实验室才配谈的数字。其实你厨房里就有一台“光速测量仪”。

\textbf{做法：}取出微波炉里的转盘（关键一步，转盘会让食物均匀受热，我们要的就是不均匀）。把一块平盘巧克力、一片芝士或一层铺在盘子里的黄油放进去，用小火加热十几秒，直到表面出现间隔均匀的软化斑或焦斑。关掉微波炉，用尺子量出相邻两个斑块中心之间的距离 \(d\)。

\textbf{你会看到：}\(d\) 大约是 \(6\,\mathrm{cm}\)。

\textbf{为什么：}微波炉里的微波在炉壁之间来回反射，形成驻波（第 35 章讲过驻波）。驻波的波腹处电场最强，食物先在那里软化；相邻波腹的距离是半个波长。所以微波的波长 \(\lambda=2d\approx 12\,\mathrm{cm}\)。炉体铭牌上写着微波频率 \(f=2450\,\mathrm{MHz}=2.45\times10^{9}\,\mathrm{Hz}\)。波速等于频率乘波长：
\[
c=f\lambda\approx 2.45\times10^{9}\times 0.12\approx 2.9\times10^{8}\,\mathrm{m/s}.
\]
和精确值 \(3.0\times10^{8}\,\mathrm{m/s}\) 只差百分之几。用一盘巧克力，你复现了物理学史上最重要的常数之一。

\textbf{想一想：}这个实验里没有任何东西在高速运动。它告诉你“真空中的光有一个确定的速度值”，但它没有回答一个更深的问题：这个数值是相对于谁而言的？相对于微波炉？相对于地球？相对于某个看不见的“空间本身”？请带着这个问题往下走。
\end{experiment}

作为对照，再想一个关于声音的经验。无风的操场上，同伴在一百米外拍手，声音大约 \(0.3\) 秒后传到你耳朵里，声速约 \(340\,\mathrm{m/s}\)。如果刮着从你背后吹向同伴的大风，你会发现声音传得“更费劲”——严格说，声波相对空气的速度是 \(340\,\mathrm{m/s}\)，而空气本身在相对地面流动，所以地面上的你会测到顺风快、逆风慢。波速是相对于“介质”的。这个经验如此自然，以至于两百年前的物理学家想都没想就断定：光，一定也是某种介质里的波。

其实你对“光比声快得多”早有身体经验。夏夜雷雨，你先看见闪电，过几秒才听见雷声。数出这段秒数，除以三，就是闪电离你大约几公里：声音三秒跑一公里，而光跑同样的距离只要三百万分之一秒，完全来不及察觉。古人用它估雷暴远近，我们今天用它体会一件事：光速和声速之间隔着六个数量级。声速的“相对介质”规律我们天天验证，光速的脾气却从没人在日常里摸到过——这正是它直到一百多年前才被查清的原因。

\step{旧模型遇到麻烦}

\textbf{第一块基石：伽利略的速度加法。}本书力学篇反复使用过一条规则：如果车厢以速度 \(v\) 相对地面运动，球以速度 \(u'\) 相对车厢运动，那么球相对地面的速度是
\[
u=u'+v.
\]
它背后是一套“换参考系”的操作：地面坐标 \(x\) 和车厢坐标 \(x'\) 之间只差一个随时间均匀增长的平移，\(x'=x-vt\)，而时间在两个参考系里被默认为同一个，\(t'=t\)。这套操作叫伽利略变换。它如此显然，以至于三百年里没有人觉得它算一个“假设”——它只是“测量”的同义词。

\textbf{第二块基石：波需要介质。}声波是空气的振动，水波是水面的振动，绳波是绳子的振动。每一种波的速度都相对于它的介质：声速相对空气是 \(340\,\mathrm{m/s}\)，你跑着去追一列声波，它相对你就会变慢。于是十九世纪的物理学家顺理成章地推断：光也是一种波（第 35 章的干涉衍射实验铁证如山），那么光波也一定有它的介质。这种假想介质被命名为“以太”——一种填满整个宇宙、绝对静止、让光波在其中传播的“空气”。光速 \(c\)，就是光相对以太的速度。

\textbf{麻烦出在第三块基石上。}第 25 章讲过，麦克斯韦把电学和磁学写进了一组方程。这组方程有一个惊人的推论：真空中，变化的电场与变化的磁场在空间和时间上相互配合，可以以波的形式自我维持地传播，传播速度由两个可以从桌面实验测出的常数决定：
\[
c=\frac{1}{\sqrt{\varepsilon_0\mu_0}}.
\]
请盯住这个式子看三秒钟。式子里有真空电容率 \(\varepsilon_0\)、真空磁导率 \(\mu_0\)，都是测量静电力和磁力时得到的常数。式子里\textbf{没有}光源的速度，\textbf{没有}观察者的速度，甚至\textbf{没有}以太的位置。麦克斯韦方程仿佛在说：电磁波的速度是一个从自然规律里直接算出来的数，跟谁在看毫无关系。

矛盾就此成型。如果伽利略变换是对的，那么同一个光速在不同参考系里应该测出不同的值，麦克斯韦方程就只能在\textbf{唯一一个}特殊参考系——以太静止系——里精确成立，搬到火车上就得出错。换句话说，电磁学似乎能分辨出“谁在绝对地运动”。可力学里的一切实验（伽利略船舱里掷球、滴水的著名论证）都说：匀速运动的人做任何力学实验都察觉不到自己在动。力学说所有惯性系平权，电磁学却指着一个特殊参考系。两块基石，必有一块要裂。

\textbf{那就让实验来裁决。}地球绕太阳公转的速度约 \(30\,\mathrm{km/s}\)。如果以太真的填满了空间并且相对太阳静止，地球就是在“以太风”里穿行：朝公转方向发出的光应该慢一点点，逆着发出的光应该快一点点，差别约为 \(v/c\approx 10^{-4}\)，即万分之一。这个差别直接测时间太难，但光的波长很短，用干涉仪可以把速度的万分之一放大成看得见的条纹移动。

一八八七年，迈克耳孙和莫雷把干涉仪做到了极致：光在臂上来回反射多次，等效臂长约 \(11\,\mathrm{m}\)。按照以太理论，把整个装置旋转九十度，干涉条纹应该移动约
\[
\Delta N=\frac{2L}{\lambda}\cdot\frac{v^{2}}{c^{2}}
=\frac{2\times11}{5.5\times10^{-7}}\times10^{-8}\approx 0.4\ \text{个条纹}.
\]
仪器能读出百分之一条纹的移动，而 \(0.4\) 是它的四十倍。结果呢？几乎为零。春天做是零，半年后地球公转方向反过来了再做，还是零。把仪器拖到高山上做，依然是零。预期的以太风，在最该出现的时候，始终不肯出现。

还有一桩更早的旧案值得一提。一八五一年，菲佐让光顺着和逆着流动的水传播，测量光速的变化。如果流水完全带动以太，岸上测得的光速应该是“光在静水中的速度”加减完整的水速；实测出来的“带动”却只有水速的约 \(0.43\) 倍。这个诡异的系数当年被解释成“以太被介质部分拖曳”，还成了以太理论一度得意的证据。请把这个 \(0.43\) 记在心里——本章后面它会自己从公式里走出来。

\begin{center}
\begin{tikzpicture}[font=\small,>={Stealth}]
  \node[draw,rectangle,inner sep=3pt] (src) at (-2.8,0) {光源};
  \draw[fill=blue!10,rotate=45] (-0.16,-0.16) rectangle (0.16,0.16);
  \node[below right=1pt] at (0.15,-0.2) {分光镜};
  \draw[very thick] (3.4,-0.5) -- (3.4,0.5) node[midway,right] {镜 1};
  \draw[very thick] (-0.5,3.4) -- (0.5,3.4) node[midway,above] {镜 2};
  \draw[-{Stealth}] (src.east) -- (-0.2,0);
  \draw[-{Stealth}] (0,0) -- (3.2,0);
  \draw[-{Stealth}] (3.2,0) -- (0.2,0);
  \draw[-{Stealth}] (0,0) -- (0,3.2);
  \draw[-{Stealth}] (0,3.2) -- (0,0.2);
  \node[draw,rectangle,inner sep=3pt] (det) at (0,-2.1) {探测器（看条纹）};
  \draw[-{Stealth}] (0,0) -- (det.north);
  \draw[-{Stealth},dashed] (-2.3,1.7) -- (2.3,1.7)
    node[midway,above] {假想的以太风};
\end{tikzpicture}
\end{center}

\textbf{救场的尝试反而越救越糟。}有人说：也许以太被地球拖着走，所以地面附近没有风。但早在一七二九年，布拉德雷就观测到恒星光行差：由于地球公转，望远镜要像“雨中打伞”一样微微倾斜，周年摆动约 \(20.5\) 角秒。如果地球附近的以太被完全拖走，这个倾斜角就不该存在——两种说法互相打架。也有人（斐兹杰惹、洛伦兹）提出：也许运动方向上的臂会恰好缩短，把条纹移动抵消掉。可这个缩短不多不少正好抵消实验结果，不为任何其他现象所必需，也不能预言任何新东西——它不是解释，是照着答案订做的补丁。一个理论走到需要为每个实验单独缝一块补丁的地步，离重写就不远了。

\step{发明新概念}

一九〇五年，在专利局当职员的爱因斯坦发表《论动体的电动力学》。他的做法出人意料：不去修改麦克斯韦方程（它们在所有实验里都对），也不去修补以太（它每次都缺席），而是回过头来，审视那个从来没人当作假设的假设——“速度当然相加”“时间对所有人一样流逝”。他把两条陈述放在整栋物理学大厦的最底层，当作新的地基：

\textbf{假设一（相对性原理）：}在所有惯性参考系中，物理规律具有相同的形式。也就是说，不只是力学实验，包括电磁学、光学在内的\textbf{任何}实验，都无法告诉你自己在做匀速直线运动。伽利略在船舱里的体验，对麦克斯韦方程同样成立。

\textbf{假设二（光速不变原理）：}真空中的光速，对一切惯性观察者都是同一个值 \(c\)，与光源的运动无关，与观察者的运动无关。

用本书每条重要公式都要回答的四个问题来检查这两条假设：

\textbf{描述什么？}任何人在任何惯性系里测量真空光速这个行为，以及物理定律在换参考系时的行为。\textbf{为什么这样写？}因为麦克斯韦方程给出的 \(c\) 不含参考系，干脆把这份“不合群”提升为普遍原理；因为以太在所有实验里探测不到，干脆承认它多余。\textbf{何时成立？}惯性系、引力可以忽略的场合。\textbf{何时失效？}引力场中、加速参考系里，需要更一般的理论——那是第 41 章的故事。

这里必须划清本书一直强调的三层界限。\textbf{实验事实}是：迄今为止一切测量都给出同一个光速，一切寻找以太风的实验都失败。\textbf{数学模型}是：从两条假设出发可以严格推出一整套新力学和新时空观（狭义相对论），它的每一个可检验预言都被证实。\textbf{物理解释}——“时间和空间到底是什么”——则允许更深的追问，后文两章会一层层剥开。把任何一层说成另外一层，都是对本章的误读。

爱因斯坦那篇论文还有一个不那么显眼、却同样锋利的动作：他把每个词都逼问成“怎么测”。他说，物理概念必须能落实到具体的测量操作，否则就是空话。比如“北京和上海的两座钟指着同一时刻”这句话，到底怎么验证？唯一的办法是交换光信号、扣除光在路上花的时间——而光在路上花多少时间，又取决于你怎么规定两地的“同时”。循环就此露头：\textbf{“同时”不是自然界盖好的章，而是一种需要用光信号约定的规矩。}沿着这个操作主义的追问再走一步，就会撞上“不同参考系对同时有不同约定”这堵墙。那堵墙是第 39 章的正题，本章只需记住方法论：凡是说不清怎么测的物理量，都先别忙着相信它存在。以太就是这样被请下台的——它不是被某个实验“证明不存在”，而是被发现根本无法测量，因而在物理学里没有立足之地。

最后，预先堵住一个流传极广的误读：狭义相对论\textbf{不是}“一切都是相对的”。恰恰相反，它的骨架是一串\textbf{绝对}的东西：光速 \(c\) 对所有人相同；后面两章会看到，事件之间的“时空间隔”对所有人相同，物体的静质量对所有人相同。真正相对化的，是过去被误认为绝对的东西——长度、时长、“同时”。这个理论叫“相对论”，但它是靠不变量站起来的。

\step{一句话模型}

\begin{oneline}
光速对每个人都一样，所以“过了多久”“相距多远”不能对每个人都一样——时间和长度必须为光速让路。
\end{oneline}

球的速度相加、光的速度不相加，区别不在球和光的“脾气”，而在于“速度等于距离除以时间”这句话里，距离和时间本身在不同参考系里丈量出不同的结果。光速不变是因，“多长”“多久”因人而异是果。这个果实怎么长出来，是第 39 章的主线；本章先把“因”钉牢。

\step{数学翻译器}

先把旧的规则写清楚，再看它在哪里崩掉。伽利略的速度加法：
\[
u=u'+v,
\]
其中 \(u'\) 是物体相对运动参考系的速度，\(v\) 是参考系相对地面的速度，\(u\) 是物体相对地面的速度。公式一出，立刻用特殊情况拷问它：\(v=0\) 时退回 \(u=u'\)，没问题；\(u'=0\) 时物体随参考系走，\(u=v\)，没问题；方向反转，\(v\) 取负号，依然自洽。可是当代入 \(u'=c\)：\(u=c+v\ne c\)——只要 \(v\ne0\)，它就和假设二迎头相撞。旧公式在光速面前，数学上就是错的。

新的速度合成规则是
\[
u=\frac{u'+v}{1+\dfrac{u'v}{c^{2}}}.
\]
同样立刻拷问它：

\textbf{低速极限：}\(u'\) 和 \(v\) 都远小于 \(c\) 时，分母里的 \(u'v/c^{2}\) 小到可以忽略，公式退回 \(u\approx u'+v\)。旧规则不是错了，而是新规则的低速近似——所以你的前半生没有被它骗过。

\textbf{代入光速：}令 \(u'=c\)，
\[
u=\frac{c+v}{1+\dfrac{cv}{c^{2}}}=\frac{c+v}{1+\dfrac{v}{c}}
=\frac{c+v}{\dfrac{c+v}{c}}=c,
\]
无论 \(v\) 是多少，结果严丝合缝地等于 \(c\)。假设二被公式自动满足：光的速度，在任何参考系合成后还是 \(c\)。

\textbf{方向反转：}令 \(u'=-c\)（光朝反方向跑），同样算出 \(u=-c\)。正负光速都是不变的。

\textbf{两个亚光速相加：}令 \(u'=0.5c\)、\(v=0.5c\)，得
\[
u=\frac{0.5c+0.5c}{1+0.25}=\frac{c}{1.25}=0.8c.
\]
旧公式会说 \(1.0c\)，新公式说 \(0.8c\)。可以证明，两个亚光速永远合不出超光速：只要 \(u'<c\) 且 \(v<c\)，结果必然小于 \(c\)。\(c\) 在公式里自动扮演“天花板”的角色。

再来一个带真实数据的估算，体会新旧公式差多少。高铁时速 \(300\,\mathrm{km/h}\)，即 \(v\approx83\,\mathrm{m/s}\)，于是 \(v/c\approx2.8\times10^{-7}\)。若在车厢里以 \(u'=10\,\mathrm{m/s}\) 抛球，新公式分母与 \(1\) 的偏离是
\[
\frac{u'v}{c^{2}}\approx\frac{10\times83}{(3\times10^{8})^{2}}\approx9\times10^{-15}.
\]
小数点后十四个零。全世界没有一台放在站台上的仪器能分辨这个差别——这就是伽利略加法统治人类直觉三百年的原因。

再看一个中等速度的例子，让新旧差别显形。设想一艘飞船相对地球以 \(0.6c\) 飞行，从飞船上向前发射一枚相对飞船 \(0.7c\) 的炮弹。旧公式预言地面看到 \(1.3c\)——超光速。新公式给出
\[
u=\frac{0.6c+0.7c}{1+0.6\times0.7}=\frac{1.3c}{1.42}\approx0.915c,
\]
比 \(c\) 小，而且永远小。分母上那个看似不起眼的修正项，正是自然界防超光速的保险丝。

还记得菲佐实验里那个 \(0.43\) 吗？现在让它回家。光在静水中的速度是 \(c/n\)（水的折射率 \(n\approx1.33\)，第 36 章），水本身以速度 \(v\) 流动。用速度合成公式把二者合成，对远小于 \(c\) 的水速做近似，得到的“被带动的比例”正好是 \(1-1/n^{2}\approx0.43\)。当年需要“以太被部分拖曳”这个专门假说才能解释的系数，如今只是速度合成公式的一阶近似。一个好理论的标志就是这样：旧理论里东拼西凑的经验系数，在新理论里变成无需假设的副产品。

在真正的加速器世界里，新旧公式已经不是修正与被修正的关系，而是对与错的关系。欧洲核子中心的对撞机让两束质子以各自约 \(0.999\,999\,99\,c\) 的速度对撞：按旧加法，相对速度接近 \(2c\)；按新公式合成，任何一束看另一束仍然小于 \(c\)。加速器工程师设计束流时序、计算对撞能量时，用的是本章这个分母，而不是小学加法——这是相对论每天都在上班的地方。

\begin{mathbridge}
\textbf{两条假设如何逼出新的时空变换。}只凭“相对性原理加光速不变”，加上变换必须是线性的（匀速运动在任何参考系看都该是匀速的），就能把伽利略变换唯一地替换成洛伦兹变换。勾画关键三步：设车厢系相对地面系以速度 \(v\) 沿 \(x\) 方向运动，线性变换最一般的形式是
\[
x'=\gamma(x-vt),\qquad t'=\gamma\!\left(t-\frac{vx}{c^{2}}\right),
\]
其中 \(\gamma\) 是待定系数。现在让一束光在 \(t=0\) 时从原点发出：地面系里它的轨迹是 \(x=ct\)，车厢系里按假设二必须是 \(x'=ct'\)。把 \(x=ct\) 代入上面两式，令 \(x'=ct'\)，整理后立刻得到
\[
\gamma=\frac{1}{\sqrt{1-v^{2}/c^{2}}}.
\]
注意两件事。第一，\(t'\) 的表达式里出现了 \(x\)：车厢系的“现在几点”，竟然取决于“在哪里”。时间的普适性就是从这里开始瓦解的——第 39 章会沿着这条裂缝拆出同时性的相对性、时间膨胀和长度收缩。第二，对上面的变换求速度 \(u=dx/dt\) 与 \(u'=dx'/dt'\) 的关系，正好得到主线里那个速度合成公式。所以它不是一个凑出来的补丁，而是两条假设的直系后代。本章不展开推导细节，你只要记住：整台机器只有两个输入——相对性原理和光速不变。
\end{mathbridge}

\begin{challenge}
\textbf{给速度换个坐标：快度。}速度合成公式之所以难看，是因为“速度”这个变量本身选得不好。定义一个叫做\textbf{快度}的量 \(\varphi\)，满足 \(v=c\tanh\varphi\)（\(\tanh\) 是双曲正切函数）。代入合成公式，利用双曲函数的加法公式 \(\tanh(\varphi_1+\varphi_2)=\dfrac{\tanh\varphi_1+\tanh\varphi_2}{1+\tanh\varphi_1\tanh\varphi_2}\)，你会惊喜地发现：\textbf{快度就是简单相加的}，\(\varphi=\varphi_1+\varphi_2\)。相对论里的“速度相加”其实是“快度相加”戴了层面具。这个变量立刻解释了两件事：为什么日常速度近似相加（\(\varphi\) 很小时 \(\tanh\varphi\approx\varphi\)，快度近似等于 \(v/c\)）；以及为什么光速永远追不上（\(v\to c\) 对应 \(\varphi\to\infty\)，无穷大快度再加任何有限值还是无穷大）。第 40 章讨论相对论能量时，快度还会再出现。
\end{challenge}

\step{模型的边界}

\textbf{适用范围。}狭义相对论的两条假设只在惯性系里、引力可以忽略时成立。它不是“终极理论”，而是一套在极其宽广的范围内从未失手的工作模型：从医院的直线加速器到宇宙线里的高能粒子，它的预言每次都对。一旦引力变得重要——黑洞附近、中子星表面，甚至你手机里的导航卫星做精密校时——就必须请出广义相对论（第 41 章）。一个常被提起的工程事实是：导航卫星上的原子钟若不按相对论（狭义加广义两部分）做修正，累积误差会让定位每天偏掉大约十公里量级，导航会在几天内报废。修正的具体机制分属第 39 章和第 41 章，这里只记住结论：你口袋里的手机，每天都在消费本章的两条假设。

\textbf{低速近似何时够用。}只要 \(v/c\) 很小，伽利略加法的误差量级是 \((v/c)^{2}\)。民航客机约 \(250\,\mathrm{m/s}\)，误差约 \(7\times10^{-13}\)；步枪子弹约 \(10^{3}\,\mathrm{m/s}\)，误差约 \(10^{-11}\)。工程上凡涉及宏观物体，牛顿力学依然是首选工具——不是因为它“更正确”，而是因为在这个尺度上它足够对，而且简单得多。判断用哪套理论，看的不是新旧，而是 \(v/c\) 的大小。

\textbf{一个悬着的证据，留给下一章。}本章只建立了两条假设和速度合成规则，还没有动用它们最惊人的预言。先埋一粒种子：宇宙射线在十几公里高空打出一种叫缪子的粒子，它的寿命只有约 \(2.2\) 微秒。按经典算法，就算它以接近光速飞行，衰变前也只能跑约 \(660\,\mathrm{m}\)，理应全部死在高空；可地面实验室每天都能抓到大把缪子。它们是怎么“活着”走完十几公里的？任何不涉及时间和长度本身的解释都救不了这个账。第 39 章会用光钟把这个悬案拆开，那时你会看到“光速不变”如何反过来改写“多久”和“多远”的含义。

\textbf{典型误用一：“相对论说明一切都是相对的。”}前面说过，正好相反：理论靠不变量站立。下次再听到有人用“爱因斯坦证明了没有绝对真理”来讨论物理学之外的话题，你可以礼貌地指出：他证明的是光速和时空间隔的绝对性。

\textbf{典型误用二：把光速不变当成“光的怪脾气”。}光只是第一个把时空结构暴露出来的信使。真正改变的是速度、长度、时间之间的关系。特别地，光在水和玻璃里传播得比 \(c\) 慢，这不违反任何东西：\(c\) 是\textbf{真空}中的极限速度，介质里的光速是光与物质相互作用的结果（第 36 章），两码事。

\textbf{典型误用三：“物体接近光速时质量变大，所以推不动。”}这是旧教材的表述方式。本书采用现代惯例：质量（静质量）是不变量，不随速度改变；随速度增长的是总能量和动量。第 40 章会用不变质量和能量—动量关系把这件事说清楚，届时你会看到“为什么加速到 \(c\) 需要无穷大的能量”有一个干净的表述，不需要“质量变大”这个拐杖。

\textbf{还有哪些“超光速”不违反本章？}扫过地面的探照灯光斑可以移动得比 \(c\) 快——光斑不是物体，不携带能量和信息从源头跑到远处；波的相速度在某些介质里可以超过 \(c\)——但相速度不传递信号。一切“超光速”的传闻，只要问一句“它能不能用来把一条消息送到比光更快的地方”，就会现出原形。二〇一一年曾轰动一时的“中微子超光速”实验，最后查明原因是一个松动的光纤接头让计时差了几十纳秒。凡声称推翻光速极限的测量，先怀疑尺子和钟——这不是顽固，而是一百多年来所有类似事件的共同结局。

\step{回头解释开场现象}

现在，回到站台。

手电在火车上朝前发光，站台上的仪器测得 \(c\)，车厢里的仪器也测得 \(c\)。为什么球要加、光不加？因为“速度等于距离除以时间”，而不同参考系里的尺子和时钟读数之间的关系，不再是伽利略变换，而是由两条假设逼出来的新变换；在新变换下，只有 \(c\) 这个速度在合成后保持不变。光速不变不是测量事故，是时空结构的线索：自然界用光告诉我们，\(c\) 是编织进空间和时间关系里的一个常数。

再看那个十六岁的问题：追着光跑会看见什么？现在你有了完整的回答。假设你坐进飞船，加速到 \(0.9c\)，再加速到 \(0.99c\)、\(0.999\,999c\)——回头看那束从地球出发的激光，它相对你的速度每一次都还是 \(c\)。你\textbf{永远}追不平它，所以“与光并肩而立、看一列冻结的波”这个画面在自然界里根本不存在。从麦克斯韦方程这边看，结论也一样：一列不传播的、冻结的电磁波，在任何参考系里都不满足方程，所以大自然从不展出这种东西，你在任何飞船上都不会看到它。两条路通向同一个答案，这正是理论自洽的样子。

至于“为什么人（以及任何有质量的东西）加速不到 \(c\)”——这需要谈能量。第 40 章会证明：越接近 \(c\)，同样的推力换来的速度增量越小，要再快一点点，所需的能量越来越大，直到无穷。\(c\) 不是一道栅栏，而是一条渐近线：你可以无限靠近，永远抵达不了。

最后清点一下旧冲突是如何化解的。麦克斯韦方程算出的光速不含参考系，曾让它看起来只能在一个特殊参考系里成立；现在，光速不变原理让麦克斯韦方程在\textbf{每一个}惯性系里以同样的形式成立。电磁学没有输，力学相对性原理也没有输——输掉的是“速度简单相加”这条从未被怀疑的常识，以及为挽留它而发明的以太。物理学的前进有时不是找到新东西，而是承认某个旧假设从来不是必需。

\step{脑洞实验与挑战题}

\begin{thought}[星系尺度的追光]
把追光实验推到极限。设想一艘飞船以 \(0.9999\,c\) 掠过地球，去追一束同向飞行的激光。地面上的人测：激光在前，飞船在后，二者的间距以 \(c-0.9999c=0.0001c\) 的速度缓慢拉大。飞船上的人测：那束激光以整整 \(c\) 的速度离船而去。两个说法差了一万倍，却\textbf{都对}——因为它们描述的是不同参考系里用各自的尺子和钟测出的不同量。“追”这个动作在日常里的全部含义（我加速，相对速度就减小），在接近光速时整个失效。再问一步：如果飞船继续加速到 \(0.999\,999\,999\,9\,c\) 呢？答案你已经知道了：那束光还是以 \(c\) 离它而去。从飞船的视角看，追光就像追自己的地平线。
\end{thought}

\begin{thought}[两台对开的飞船]
两艘飞船在远离地球的深空相向而行，地面测得它们各自的速度都是 \(0.75c\)。地面上的人会算出：两船以 \(1.5c\) 的“接近速率”互相靠近。这违反光速极限吗？不违反。光速极限约束的是：任何一个观察者测得的、某个物体相对自己的速度。地面上的你并没有看到任何一个物体相对你超光速；而坐在任一艘飞船里的人，用本章的速度合成公式去测另一艘船，得到的是
\[
u=\frac{0.75c+0.75c}{1+0.75\times0.75}=\frac{1.5c}{1.5625}=0.96c,
\]
依然小于 \(c\)。“第三者眼中的相互接近速率”和“当事人眼中的相对速度”是两个不同的量，前者可以超过 \(c\) 而不传递任何信息，后者永远被压在 \(c\) 以下。分清这两个量，是读科幻新闻时的防身术。
\end{thought}

\begin{puzzles}
\item 火车上的人向前抛出 \(10\,\mathrm{m/s}\) 的球，站台测得球速是车速加球速；他改用激光笔朝前照，站台测得的光速却不加。用本章的两条假设说明：差别究竟出在哪里？（思路提示：不要问“光为什么特殊”，而要问“速度相加这条规则是哪来的”。它来自伽利略变换，而伽利略变换只是低速近似。）
\item 假如以太真的存在，而且在地球附近被地球完全拖着走。这个“拖曳说”能解释迈克耳孙—莫雷实验的零结果，却会在恒星光行差（望远镜因地球公转而需倾斜约 \(20.5\) 角秒）面前栽跟头。说明为什么。（思路提示：被地球拖着一起走的以太，会让远处射来的星光在进入“以太拖曳区”时发生偏折，把光行差抹掉或改变它的大小；但光行差早在一七二九年就被观测到了。）
\item 用速度合成公式算三笔账：(a) \(u'=0.5c,\ v=0.5c\)；(b) \(u'=c,\ v=0.9c\)；(c) \(u'=10\,\mathrm{m/s},\ v=100\,\mathrm{km/h}\)。（思路提示：(a) 答案不是 \(c\)；(b) 答案必须是 \(c\)，算出来不是 \(c\) 就说明代错了；(c) 先估一下分母与 \(1\) 的偏离量级，你会发现直接用小学加法完全合法——学会判断“什么时候可以偷懒”也是物理能力。）
\item 有人这样解释光速不变：“运动的尺子收缩、运动的钟变慢，所以仪器被‘骗’了，量出来的光速才总是 \(c\)。”这个说法漏掉了什么？（思路提示：如果是仪器被骗，不同原理的钟——机械钟、原子钟、人的心跳——没有理由被骗得一模一样；而且“骗局”必须在任何速度、任何方向上恰好给出同一个 \(c\)。第 39 章会看到，真相更大胆：不是仪器被骗，而是“多长”“多久”本身确实依赖参考系。）
\end{puzzles}
